\documentclass[12pt]{article}
\usepackage{amsmath,amssymb,amsfonts,amsthm}
\usepackage{mathrsfs}
\usepackage{enumitem}
\usepackage[colorlinks=true,linkcolor=blue,citecolor=blue,urlcolor=blue]{hyperref}

\textheight=21.4cm
\textwidth=15cm
\topmargin=-0.2in
\oddsidemargin=0cm
\evensidemargin=-0.2cm

\newtheorem{theorem}{Theorem}[section]
\newtheorem{proposition}[theorem]{Proposition}
\newtheorem{lemma}[theorem]{Lemma}
\newtheorem{corollary}[theorem]{Corollary}
\newtheorem{definition}[theorem]{Definition}
\newtheorem{problem}[theorem]{Problem}
\newtheorem{remark}[theorem]{Remark}

\newcommand{\A}{\mathcal A}
\newcommand{\B}{\mathcal B}
\newcommand{\J}{\mathcal J}
\newcommand{\K}{\mathcal K}
\newcommand{\CB}{\mathrm{CB}}
\newcommand{\id}{\mathrm{id}}
\newcommand{\otmin}{\otimes_{\min}}
\newcommand{\otmax}{\otimes_{\max}}
\newcommand{\Ball}{\mathrm{Ball}}
\newcommand{\dex}{d_{\mathrm{ex}}}
\newcommand{\dcb}{d_{\mathrm{cb}}}
\newcommand{\Mn}{M_n}

\title{Local Reflexivity versus Exactness: Effros--Ruan Conditions, Stability of the Pisier Constant, and an Open Problem}
\author{Yafei Zhao}
\date{}

\begin{document}
\maketitle

\begin{center}
\begin{minipage}{11.5cm}
{\bf Abstract.}
Kirchberg proved that every exact $C^*$-algebra is locally reflexive; the converse is open.  This note does not settle that converse.  Its mathematical content is twofold.  First, we record an elementary but useful \emph{stability lemma}: the Pisier exactness constant $\dex$ is controlled under completely bounded isomorphisms with explicit constants.  As a consequence, any finite-dimensional subspace of a non-exact $C^*$-algebra carries a quantitative obstruction that no local reflexivity map can erase.  Second, we foreground the Effros--Ruan \emph{tensor--bidual} conditions $C'$ and $C''$: exactness is equivalent to $C'$ (for every operator space $W$, the canonical map $\A\otmin W^{**}\to(\A\otmin W)^{**}$ is a complete isometry), while local reflexivity is equivalent to $C''$ (the analogous map with $\A^{**}$ in place of $W^{**}$).  Together with the lemma ``$C$ if and only if $C'$ and $C''$'', this packages the open implication as ``does $C''$ alone force $C'$?''---logically the same strength as exactness versus local reflexivity, but with a narrative centred on minimal tensor products and biduals rather than on choosing images of maps into $\A$.  We append a short remark on an equivalent \emph{matricial} restatement in terms of $\dex$ and Lemma~\ref{lem:stab}.  We close with proved sufficient conditions for exactness and a conservative comparison of local reflexivity with the local lifting property.

\medskip
{\bf Keywords.} exact $C^*$-algebra, local reflexivity, condition $C'$, condition $C''$, Pisier constant, completely bounded map, operator space.

\medskip
{\bf 2020 Mathematics Subject Classification.} Primary 46L06; Secondary 46L07, 46M05.
\end{minipage}
\end{center}

\section{Introduction}

Write $\A$ for a $C^*$-algebra.  \emph{Exactness} is the tensorial preservation of short exact sequences under $\A\otmin -$.  \emph{Local reflexivity} is a bidual finite-dimensional approximation principle: finite-dimensional subspaces of $\A^{**}$ admit linear maps into $\A$ that are nearly complete isometries on the subspace and interpolate finitely many prescribed values on $\A^*$.  Kirchberg proved
\[
        \A\ \text{exact}\quad\Longrightarrow\quad \A\ \text{locally reflexive};
\]
see \cite[Ch.~9]{BO}.  The reverse implication is a long-standing open problem.

Effros and Ruan reformulated both properties as \emph{isometric tensor--bidual identities} valid for all operator spaces $W$: exactness is their condition $C'$ (tensor with the bidual of $W$ on the right), while local reflexivity is condition $C''$ (tensor with $\A^{**}$ on the left).  Condition $C$ is the conjunction $C'\wedge C''$ and is equivalent to a natural injective-tensor identity; see \cite[\S17.1, Thm.~PEX.1, Thm.~PLR1.1]{ER} (together with the lemma recorded below as Lemma~\ref{lem:ccc}).  For $C^*$-algebras, Kirchberg's theorem implies that exactness already forces $C''$, so $C'$ and $C''$ hold together whenever $\A$ is exact.  The unresolved converse ``locally reflexive$\Rightarrow$exact'' is thus equivalent to asking whether $C''$ alone implies $C'$ for $C^*$-algebras.

The stability lemma in \S\ref{sec:stab} is independent of that dictionary: it isolates a finite-dimensional mechanism for why non-exact algebras carry obstructions inside $\A$ itself.  \S\ref{sec:cc} states $C'$ and $C''$ and records the cited equivalences without reproving the deep characterizations.

\section{Preliminaries}\label{sec:prelim}

We use operator spaces in the sense of Effros and Ruan \cite{ER}.  For operator spaces $V$ and $W$,
\[
        \dcb(V,W)=\inf\bigl\{\|u\|_{\mathrm{cb}}\,\|u^{-1}\|_{\mathrm{cb}} : u:V\to W\ \text{a completely bounded isomorphism}\bigr\}.
\]
For finite-dimensional $L\subseteq V$ define
\begin{align*}
        \dex(L) &= \inf\bigl\{\dcb(L,S): n\in\mathbb N,\ S\subseteq \Mn\bigr\},\\
        \dex(V) &= \sup\bigl\{\dex(L): L\subseteq V\ \text{finite-dimensional}\bigr\}.
\end{align*}
By Pisier's theorem \cite[Thm.~17.1]{Pisier} (see also \cite[Thm.~14.4.1]{ER}), a $C^*$-algebra $\A$ is exact if and only if $\dex(\A)=1$.

\begin{definition}[Local reflexivity]\label{def:lr}
The $C^*$-algebra $\A$ is \emph{locally reflexive} if for every finite-dimensional $E\subseteq \A^{**}$, every finite-dimensional $F\subseteq \A^*$, and every $\varepsilon>0$, there exists a linear map $\phi:E\to \A$ such that
\begin{enumerate}[label=(\roman*),leftmargin=2.4em]
\item $\|\phi\|_{\mathrm{cb}}\,\|(\phi_{|\phi(E)})^{-1}\|_{\mathrm{cb}} < 1+\varepsilon$,
\item $f(\phi(x))=x(f)$ for all $x\in E$ and $f\in F$.
\end{enumerate}
\end{definition}

\begin{definition}[Exactness]
The $C^*$-algebra $\A$ is \emph{exact} if for every short exact sequence
$0\to \J\to \B\overset{\pi}{\to}\B/\J\to 0$ of $C^*$-algebras, the induced sequence
$0\to \A\otmin\J\to \A\otmin\B\overset{\id\otimes\pi}{\longrightarrow}\A\otmin(\B/\J)\to 0$
is exact in the operator-space sense.
\end{definition}

For a closed ideal $\J\triangleleft \B$, the \emph{Fubini product} is
$F(\A,\B,\J)=\{u\in \A\otmin\B : (f\otimes\id_\B)(u)\in \J\ \text{for all}\ f\in \A^*\}$.
As in Lemma~2.4 of an earlier draft, $F(\A,\B,\J)$ agrees with the kernel of $\id_\A\otimes\pi$.

\begin{theorem}[Effros--Ruan]\label{thm:plr14}
If $\A$ is a locally reflexive $C^*$-algebra and $\J\triangleleft\A$, then for every operator space $W$,
\[
        \ker\bigl(\pi\otimes\id_W:\A\otmin W\to (\A/\J)\otmin W\bigr)=\J\otmin W.
\]
\end{theorem}

\begin{proof}[Reference]
See \cite[\S17]{ER} (local reflexivity implies good behaviour of $\otmin$ with quotients when the ideal lies in the left factor); compare also \cite[Ch.~9]{BO}.
\end{proof}

\section{Stability of \texorpdfstring{$\dex$}{dex} under cb isomorphisms}\label{sec:stab}

\begin{lemma}[Stability]\label{lem:stab}
Let $L$ and $M$ be finite-dimensional operator spaces and let $T:L\to M$ be a linear isomorphism onto its image.  Then
\[
        \dex(M)\le \|T\|_{\mathrm{cb}}\,\|T^{-1}\|_{\mathrm{cb}}\,\dex(L),
        \qquad
        \dex(L)\le \|T\|_{\mathrm{cb}}\,\|T^{-1}\|_{\mathrm{cb}}\,\dex(M).
\]
In particular, if $\|T\|_{\mathrm{cb}}\,\|T^{-1}\|_{\mathrm{cb}}\le \lambda$, then
\[
        \dex(L)/\lambda \le \dex(M)\le \lambda\,\dex(L).
\]
\end{lemma}

\begin{proof}
Fix $\eta>0$ and choose $n\in\mathbb N$ and $S\subseteq \Mn$ with a linear isomorphism $u:L\to S$ satisfying $\|u\|_{\mathrm{cb}}\|u^{-1}\|_{\mathrm{cb}}<\dex(L)+\eta$.  Set $v=u\circ T^{-1}:T(L)=M\to S$.  Then
$\|v\|_{\mathrm{cb}}\le \|u\|_{\mathrm{cb}}\|T^{-1}\|_{\mathrm{cb}}$ and
$\|v^{-1}\|_{\mathrm{cb}}\le \|T\|_{\mathrm{cb}}\|u^{-1}\|_{\mathrm{cb}}$, whence
$\|v\|_{\mathrm{cb}}\|v^{-1}\|_{\mathrm{cb}}\le \|T\|_{\mathrm{cb}}\|T^{-1}\|_{\mathrm{cb}}(\dex(L)+\eta)$.
Thus $\dex(M)\le \|T\|_{\mathrm{cb}}\|T^{-1}\|_{\mathrm{cb}}(\dex(L)+\eta)$.  Let $\eta\to 0$.

For the reverse inequality apply the first part to $T^{-1}:M\to L$.
\end{proof}

\begin{corollary}[Obstruction on local reflexivity images]\label{cor:obstr}
Let $\A$ be a $C^*$-algebra.  Suppose there exists a finite-dimensional subspace $L\subseteq \A$ and a number $\delta>0$ such that $\dex(L)\ge 1+\delta$.  Fix $\varepsilon\in(0,1)$ small enough that $(1+\delta)/(1+\varepsilon)^2>1+\varepsilon$.  Then there is \emph{no} linear map $\phi:L\to \A$ satisfying simultaneously
\begin{enumerate}[label=(\alph*),leftmargin=2.4em]
\item $\|\phi\|_{\mathrm{cb}}\,\|(\phi_{|\phi(L)})^{-1}\|_{\mathrm{cb}}<1+\varepsilon$,
\item $\dex(\phi(L))\le 1+\varepsilon$.
\end{enumerate}
\end{corollary}

\begin{proof}
Assume such a $\phi$ exists.  Lemma~\ref{lem:stab} applied to $T=\phi$ gives
$\dex(\phi(L))\ge \dex(L)/(1+\varepsilon)^2\ge (1+\delta)/(1+\varepsilon)^2>1+\varepsilon$, contradicting~(b).
\end{proof}

\begin{remark}
Corollary~\ref{cor:obstr} is independent of weak$^*$ interpolation: it is a pure finite-dimensional operator-space obstruction.  It says that a non-exact algebra already carries, inside $\A$ itself, finite-dimensional data whose Pisier constant cannot be lowered below a threshold by any map that is merely a near cb-isomorphism on that subspace.
\end{remark}

\section{\texorpdfstring{Conditions $C'$ and $C''$}{Conditions C' and C''}}\label{sec:cc}

We follow Effros--Ruan \cite[\S17.1]{ER}.  For each operator space $W$, the algebraic tensor product $\A\odot W^{**}$ sits canonically in both $\A\otmin W^{**}$ and $(\A\otmin W)^{**}$; the inclusion $(\A\otmin W)^*\hookrightarrow(\A\odot W)^*$ induces a completely contractive map from $\A\otmin W^{**}$ into $(\A\otmin W)^{**}$.  Condition $C'$ asks that this map be completely isometric.  Similarly, $\A^{**}\odot W$ maps completely contractively into $(\A\otmin W)^{**}$; condition $C''$ asks for a complete isometry.

\begin{definition}[Condition $C'$]\label{def:cp}
The $C^*$-algebra $\A$ satisfies \emph{condition $C'$} if for every operator space $W$, the canonical completely contractive map
\[
        j'_W:\ \A\otmin W^{**}\longrightarrow (\A\otmin W)^{**}
\]
is a complete isometry.
\end{definition}

\begin{definition}[Condition $C''$]\label{def:cpp}
The $C^*$-algebra $\A$ satisfies \emph{condition $C''$} if for every operator space $W$, the canonical completely contractive map
\[
        j''_W:\ \A^{**}\otmin W\longrightarrow (\A\otmin W)^{**}
\]
is a complete isometry.
\end{definition}

\begin{theorem}[Effros--Ruan {\cite[Thm.~PEX.1]{ER}}]\label{thm:pex}
For a $C^*$-algebra $\A$, the following are equivalent.
\begin{enumerate}[label=(\roman*),leftmargin=2.4em]
\item $\A$ is exact;
\item $\A$ satisfies condition $C'$;
\item $\dex(\A)=1$ \textup{(Pisier).}
\end{enumerate}
\end{theorem}

\begin{theorem}[Effros--Ruan {\cite[Thm.~PLR1.1]{ER}}]\label{thm:plr}
For a $C^*$-algebra $\A$, the following are equivalent.
\begin{enumerate}[label=(\roman*),leftmargin=2.4em]
\item $\A$ is locally reflexive;
\item $\A$ satisfies condition $C''$.
\end{enumerate}
\end{theorem}

\begin{lemma}[Effros--Ruan]\label{lem:ccc}
An operator space satisfies Effros--Ruan condition $C$ if and only if it satisfies both $C'$ and $C''$.
\end{lemma}

\begin{proof}[Reference]
This is the lemma stated in \cite{ER} immediately after the definitions of conditions $C$, $C'$, and $C''$ in \S17.1.
\end{proof}

\begin{remark}[Same strength, different book-keeping]
Theorem~\ref{thm:pex} says that $C'$ is \emph{equivalent} to exactness, not a weakening of it.  The gain for exposition is purely linguistic: one may discuss exactness through tensor products and biduals rather than through short exact sequences or through matricial models for $\dex$.
\end{remark}

\begin{remark}[Matricial packaging]\label{rem:mlr}
Combining Theorem~\ref{thm:pex} with Kirchberg's theorem, if $\A$ is exact then every local reflexivity map has finite-dimensional image with $\dex=1$, hence that image is almost completely isometric to a subspace of some $\Mn$ (Pisier).  Conversely, if one could force such matricial control for \emph{every} bidual subspace $E\subseteq\A^{**}$ and every weak$^*$ interpolation constraint, then Corollary~\ref{cor:obstr} would imply exactness.  This is the same mathematics as $C'$, only phrased in terms of $\phi(E)\subseteq\A$ and $\dex$ instead of $j'_W$.
\end{remark}

\section{Proved sufficient conditions and a conservative comparison}

\subsection{Sufficient conditions}

The following statements are classical; we record them because they give \emph{proved} instances in which the open implication ``locally reflexive$\Rightarrow$exact'' holds for trivial reasons.

\begin{proposition}\label{prop:nuclear}
If $\A$ is nuclear, then $\A$ is exact and locally reflexive.
\end{proposition}

\begin{proof}
Nuclearity implies exactness \cite[Ch.~2]{BO}.  Exactness implies local reflexivity by Kirchberg's theorem.
\end{proof}

\begin{proposition}\label{prop:af}
If $\A$ is an AF $C^*$-algebra, then $\A$ is nuclear, hence exact and locally reflexive.
\end{proposition}

\begin{proof}
An AF algebra is a direct limit of finite-dimensional $C^*$-algebras, each of which is nuclear; nuclearity passes to inductive limits \cite{BO}.
\end{proof}

\subsection{Local reflexivity versus the local lifting property}

Recall that a unital $C^*$-algebra $\A$ has the \emph{local lifting property} (LLP) if for every unital inclusion $\J\triangleleft\B$ and every unital completely positive map $\theta:\A\to \B/\J$, each finite-dimensional operator subsystem of $\A$ admits a completely positive lift into $\B$ after restriction \cite[Def.~13.1.1]{BO}.  Local reflexivity and the LLP are different axioms: neither implies the other in full generality, and neither is known to imply exactness without additional hypotheses.  We do not attempt a complete diagram here; see \cite[Ch.~9, Ch.~13]{BO}.

\section{Slice kernels under local reflexivity and a reduction}

\begin{proposition}[Slice characterization]\label{prop:slice}
The algebra $\A$ is exact if and only if $F(\A,\B,\J)=\A\otmin\J$ for all $\B$ and all closed ideals $\J\triangleleft\B$.
\end{proposition}

\begin{proof}
Same as before: $F(\A,\B,\J)$ is the kernel of $\id_\A\otimes\pi$.
\end{proof}

\begin{problem}\label{prob:main}
If $\A$ is locally reflexive, must $\A$ be exact?
\end{problem}

By Theorem~\ref{thm:plr14}, local reflexivity already controls the \emph{kernel} of $\pi\otimes\id_W$ when the quotient lives inside $\A$ itself.  The content of Problem~\ref{prob:main} is therefore not kernel identification but the quotient-norm geometry of minimal tensor products with arbitrary $\B$.  In Effros--Ruan language, one asks whether condition $C''$ alone implies condition $C'$ for $C^*$-algebras.

\section*{Acknowledgements}
The stability lemma is standard operator-space bookkeeping; it is included because it isolates the finite-dimensional quantity behind Corollary~\ref{cor:obstr} and the matricial remark in \S\ref{sec:cc}.

\begin{thebibliography}{99}

\bibitem{BO}
N.~P. Brown and N.~Ozawa,
\emph{$C^*$-Algebras and Finite-Dimensional Approximations},
Graduate Studies in Mathematics, vol.~88,
American Mathematical Society, Providence, RI, 2008.

\bibitem{ER}
E.~G. Effros and Z.-J. Ruan,
\emph{Operator Spaces},
London Mathematical Society Monographs (New Series), vol.~23,
Oxford University Press, New York, 2000.

\bibitem{Pisier}
G. Pisier,
\emph{Introduction to Operator Space Theory},
London Mathematical Society Lecture Note Series, vol.~294,
Cambridge University Press, Cambridge, 2003.

\end{thebibliography}

\end{document}
